\origpage{228}{ROBERTSON AND SEYMOUR}
\noindent it is \textit{finite}. One possible explanation for this is that no one tries to find excluded minor theorems for ugly properties, and that nice properties might be expected to produce nice sets of excluded minors. But there is another possible explanation, contained in the following conjecture of K. Wagner (unpublished), which we hope to prove in a future paper.

\stmt
(1.1) Conjecture. \textit{If $G_1, G_2, \dots$ is any countable sequence of graphs then there exist $j > i \ge 1$ such that $G_i$ is isomorphic to a minor of $G_j$.}

This is connected with excluded minor theorems as follows. Let $P$ be any property of graphs which can in principle be characterized by excluded minors; that is, any graph isomorphic to a minor of a graph with property $P$ also has property $P$. Let $L'$ be the class of all minor-minimal graphs not possessing property $P$, and let $L \subseteq L'$ contain exactly one representative of each isomorphism class of $L'$. Then we can assert ``a graph has property $P$ if and only if it has no minor isomorphic to a graph in $L$.'' But no member of $L$ is isomorphic to a minor of another, and so Wagner's conjecture implies that $L$ is finite.

In this paper we prove a special case of Wagner's conjecture, that the conjecture holds for a sequence $G_1, G_2, \dots$ when $G_1$ is planar. We also prove some variations and extensions, which will be important in future papers of this series. (At the time of writing, we believe that we have a proof of (1.1) itself, and indeed of a similar conjecture, due to Nash--Williams, the ``immersion'' conjecture, although the latter has not yet been carefully written out. These proofs in particular rely heavily on the results of the present paper.)

A \textit{quasi-order} $\Omega$ is a pair $(E(\Omega), \le)$ where $E(\Omega)$ is a class and $\le$ is a binary relation on $E(\Omega)$ which is reflexive and transitive. (This would become a partial order if we made the third requirement of antisymmetry.) We shall study the quasi-orders on classes of graphs defined by the relation ``is isomorphic to a minor of.'' We treat this as a quasi-order rather than as a partial order, because it will be necessary to distinguish between isomorphism and equality. A quasi-order $\Omega$ is a \textit{well-quasi-order} if for every countable sequence $x_1, x_2, \dots$ of elements of $E(\Omega)$, there exist $i' > i \ge 1$ such that $x_i \le x_{i'}$.

The reader may perhaps consider it more natural to study quasi-orders with no infinite antichain to approach the finiteness phenomenon for excluded minor theorems, rather than to study well-quasi-orders. However, it comes to the same thing; for if $\Omega$ has no infinite descending chain (which is obvious in our context) then it is a well-quasi-order if and only if it has no infinite antichain, as may easily be verified. We choose the sequence formulation because we find it easier to manipulate.

Incidentally, there are some other natural quasi-orders of the class of all
